\uuid{VLI6}
\titre{Primitives usuelles/fractions rationnelles}
\auteur{hanani ; mzari}
\organisation{exo7}
\licence{CC BY-NC-SA 4.0}
\datecreate{2020-01-20}
\niveau{L1}
\module{Analyse}
\chapitre{Calcul d'intégrales}
\sousChapitre{Primitives diverses}
\difficulte{3}
\theme{Primitives, primitives usuelles, fractions rationnelles}
\exo7qcmid{lille-2.primitives.043}

\contenu{
\begin{question}
\qclassification{127.11}
\qtags{motcle=primitives usuelles/fractions rationnelles}
Le changement de variable $t=x^2$, pour $x\in ]0,+\infty[$, donne :
\begin{answers}
\bad{$\displaystyle \int \frac{2\mathrm{d}x}{x(x^2+1)^2}=\int \frac{\mathrm{d}t}{(t+1)^2}$.}
\good{$\displaystyle \int \frac{2\mathrm{d}x}{x(x^2+1)^2}=\int \frac{\mathrm{d}t}{t(t+1)^2}$.}
\bad{$\displaystyle \int \frac{2\mathrm{d}x}{x(x^2+1)^2}=\int \frac{\mathrm{d}t}{\sqrt{t}(t+1)^2}$.}
\good{$\displaystyle \int \frac{2\mathrm{d}x}{x(x^2+1)^2}=\int \frac{\mathrm{d}t}{t}-\int \frac{\mathrm{d}t}{t+1}-\int \frac{\mathrm{d}t}{(t+1)^2}$.}
\end{answers}
\vskip2mm
\begin{explanations}
Avec $t= x^2$, on a : $\mathrm{d}t=2x\mathrm{d}x$ et
$$\int \frac{2\mathrm{d}x}{x(x^2+1)^2}=\int \frac{2x\mathrm{d}x}{x^2(x^2+1)^2}=\int \frac{\mathrm{d}t}{t(t+1)^2}.$$
Ensuite, on décompose en éléments simples donne :
$$\frac{1}{t(t+1)^2}=\frac{1}{t}-\frac{1}{t+1}-\frac{1}{(t+1)^2}.$$
\end{explanations}
\end{question}
}
